% ! TeX root = distributed-systems-final-report.tex
\section{Concept}\label{concept}

The developed product is a simulation for an hypothetical (and simplified)
scenario where some coordination is required to guide a certain number of
autonomous cars through an intersection with a roundabout: a Cyber-Physical
System. Each vehicle is able to proceed on its own, but a central controller is
present to communicate with them, in order to optimize the traffic. The
simulation can be visualized through a web application and exposes, through a
web service, some commands, for example accessible via terminal scripts or via
the application itself.

The simulation has been thought to run endlessly without the need for human
interaction; the commands may be used for debugging or for testing particular
scenarios (pause/resume, change the state of some vehicles to disconnected from
the network, fetch information on them). No user data is thus required to be
stored.

Even inside the simulated world, cars are completely autonomous and don't
require any human interaction. The central controller only needs to collect from
each one some spatial information (such as position an speed) and store it for a
small amount of time (e.g. 1 second, after which data is considered stale and no
longer useful).

This scenario could already be considered a distributed system due to its own
nature, where the cars are nodes that, in order to navigate safely, need to at
least gather information on each other's position. Adding a central controller,
then, opens up for more possibilities to improve the efficiency of the system,
while making sure that it's safe for its participants at all times: knowing each
one's position and destination and being able to adjust the speed of each one,
it can optimize their interleaving order and make them run under more
sustainable conditions. For example, braking too hard, with time, wears some
components such as the tyres; on the other hand, avoiding higher speeds when not
need will also save some fuel, besides carrying fewer risks in general.
